% =====================================================================
%\renewcommand{\refname}{}
\begin{thebibliography}{30}

% === ITU-R Recommendations (Propagation) ===
\bibitem{itu_m2101}
ITU-R Recommendation M.2101-0, ``Characteristics of High Altitude Platform Stations
(HAPS) operating in the stratosphere,'' 2017.
Available: \url{https://www.itu.int/rec/R-REC-M.2101/latest/}

\bibitem{itu_m2106}
ITU-R Recommendation M.2106-1, ``Characteristics of High Altitude Platform Stations
(HAPS) operating in the stratosphere,'' 2012.
Available: \url{https://www.itu.int/rec/R-REC-M.2106/latest/}

\bibitem{itu_p530}
ITU-R Recommendation P.530-17, ``Propagation data and prediction methods required
for the design of terrestrial line-of-sight systems,'' 2021.
Available: \url{https://www.itu.int/rec/R-REC-P.530/latest/}

\bibitem{itu_p618}
ITU-R Recommendation P.618-13, ``Propagation data and prediction methods for the
design of Earth-space telecommunication systems,'' 2021.
Available: \url{https://www.itu.int/dms_pubrec/itu-r/rec/p/R-REC-P.618-13-202108-S!!PDF-E.pdf}
[Covers rain attenuation reduction factors, slant-path geometry, scintillation model]

\bibitem{itu_p838}
ITU-R Recommendation P.838-3, ``Specific attenuation model for rain for use in
prediction methods,'' 2005.
Available: \url{https://www.itu.int/dms_pubrec/itu-r/rec/p/R-REC-P.838-3-200503-S!!PDF-E.pdf}
[Annex 1: frequency-dependent rain coefficients $k(f)$, $\alpha(f)$]

\bibitem{itu_p676}
ITU-R Recommendation P.676-12, ``Attenuation by atmospheric gases and related
effects,'' 2021.
Available: \url{https://www.itu.int/dms_pubrec/itu-r/rec/p/R-REC-P.676-12-202112-S!!PDF-E.pdf}
[Section 3.1.2: O$_2$ resonance (59.8 GHz, 118.75 GHz); Section 3.2.1: H$_2$O resonance (22.235 GHz)]

\bibitem{itu_p840}
ITU-R Recommendation P.840-6, ``Attenuation due to clouds and fog,'' 2021.
Available: \url{https://www.itu.int/dms_pubrec/itu-r/rec/p/R-REC-P.840-6-202101-I!!PDF-E.pdf}
[Liquid water content model for fog/cloud attenuation]

\bibitem{itu_p835}
ITU-R Recommendation P.835-7, ``Reference standard atmospheres,'' 2017.
Available: \url{https://www.itu.int/dms_pubrec/itu-r/rec/p/R-REC-P.835-7-201703-I!!PDF-E.pdf}

\bibitem{itu_p836}
ITU-R Recommendation P.836-5, ``Surface water vapour statistics,'' 2017.
Available: \url{https://www.itu.int/dms_pubrec/itu-r/rec/p/R-REC-P.836-5-201703-I!!PDF-E.pdf}

\bibitem{itu_p837}
ITU-R Recommendation P.837-7, ``Characteristics of precipitation for propagation
modelling,'' 2017.
Available: \url{https://www.itu.int/dms_pubrec/itu-r/rec/p/R-REC-P.837-7-201706-I!!PDF-E.pdf}

\bibitem{itu_p1410}
ITU-R Recommendation P.1410-5, ``Propagation data and prediction methods for the
design of terrestrial broadband radioaccess systems operating in a frequency range
from 3 to 60 GHz,'' 2021.
Available: \url{https://www.itu.int/rec/R-REC-P.1410/latest/}
[Elevation-angle-dependent LoS probability baseline]

% === 3GPP / NTN Standards ===
\bibitem{3gpp_ts38821}
3GPP TS 38.821, ``Study on New Radio (NR) and its Evolution for Non-Terrestrial
Networks (NTN),'' Release 17, September 2022.
Available: \url{https://www.3gpp.org/DynaReport/38821.htm}
[HAPS frequency allocations: 38--39.5 GHz feeder, 2.1--2.2 GHz service; UAV relay scenarios]

\bibitem{3gpp_ts38214}
3GPP TS 38.214, ``NR; Physical layer procedures for data,'' Release 16, 2021.
Available: \url{https://www.3gpp.org/DynaReport/38214.htm}
[CQI/MCS-to-SINR design reference: CQI index 1 (QPSK, code rate 78/1024)
requires $\mathrm{SINR}\approx-6$\,dB; CQI index 7 (16-QAM, code rate
378/1024) requires $\mathrm{SINR}\approx+5$\,dB; CQI index 15 (256-QAM, code
rate 948/1024) requires $\mathrm{SINR}>+22$\,dB, all at the standard 10\% first-
transmission BLER target.]

\bibitem{karaman2025haps}
B.~Karaman, I.~Basturk, F.~Kara, E.~Zeydan, E.~A.~Beyazit, S.~Taskin,
E.~Bj\"ornson, and H.~Yanikomeroglu, ``On-Demand HAPS-Assisted Communication
System for Public Safety in Emergency and Disaster Response,'' \emph{IEEE
Commun. Mag.}, 2025. arXiv:2507.09153.
[Ka-band (30\,GHz) HAPS backhaul link budget, Table I: Tx power 43.2\,dBm, Tx
antenna gain 13.8\,dBi, Rx (VSAT) antenna gain 39.7\,dBi -- used here as the
literature-typical antenna-gain basis for this document's 38\,GHz
feeder-style links, \S\ref{sec:relay-path-analysis}.]

\bibitem{anim2021uav}
K.~Anim, J.-N.~Lee, and Y.-B.~Jung, ``High-Gain Millimeter-Wave Patch Array
Antenna for Unmanned Aerial Vehicle Application,'' \emph{Sensors}, vol.~21,
no.~11, p.~3914, 2021. DOI: 10.3390/s21113914
[UAV-mountable $32\times32$ mmWave patch array antenna, realized gain
30.7--32.8\,dBi at 25.4--26.9\,GHz -- basis for this document's UAV-side
receive antenna gain on the HAPS\,$\rightarrow$\,UAV hop,
\S\ref{sec:relay-path-analysis}.]

\bibitem{3gpp_ts36104}
3GPP TS 36.104, ``Evolved Universal Terrestrial Radio Access (E-UTRA); Base Station
(BS) Radio Transmission and Reception,'' Release 16, December 2021.
Available: \url{https://www.3gpp.org/DynaReport/36104.htm}
[LTE channel bandwidth specifications; 20 MHz service-link allocation]

% === Peer-Reviewed Publications ===
\bibitem{arani2023}
A.~H. Arani, P.~Hu, and Y.~Zhu, ``HAPS-UAV-Enabled Heterogeneous Networks: A Deep
Reinforcement Learning Approach,'' \emph{IEEE Open J. Commun. Soc.}, vol.~4,
pp.~1745--1760, 2023.
DOI: 10.1109/OJCOMS.2023.3294598

\bibitem{holispechac2008}
J.~Holis and P.~Pechac, ``Elevation Dependent Shadowing Model for Mobile
Communications via High Altitude Platforms in Built-up Areas,'' \emph{IEEE
Trans. Antennas Propag.}, vol.~56, no.~4, pp.~1078--1084, Apr. 2008.
DOI: 10.1109/TAP.2007.914935
[Building-blockage LoS probability series; environment parameters (Table I): $(\alpha, \beta, \xi)$ for suburban/urban/dense-urban/high-rise]

\bibitem{bakambekova2024}
A.~Bakambekova, N.~Kouzayha, and T.~Al-Naffouri, ``On the Interplay of Artificial
Intelligence and Space-Air-Ground Integrated Networks: A Survey,'' \emph{IEEE Open
J. Commun. Soc.}, vol.~5, pp.~4613--4673, 2024.

% === Regulatory / Industry Standards ===
\bibitem{ecc_report_156}
European Communications Committee (ECC), ``Minimum performance characteristics
for spectrum sharing between Fixed Satellite Service (FSS) and other services,''
ECC Report 156, April 2011 (updated December 2014).
Available: \url{https://www.ecc.cept.org/ecc/documents/ecc-reports/}
[Spectrum-sharing guidance; rain attenuation margins at Ka-band]

\bibitem{ecc_rec_04_05}
European Communications Committee (ECC), ``Radio Frequency Channel Arrangements for
Fixed Service Systems Operating in the Bands 6--7 GHz and 36--40 GHz,''
ECC Recommendation (04)05, Strasbourg, France.
[Ka-band channelization for backhaul systems]

% === SINR and Signal Analysis ===
\bibitem{itu_p618}
ITU-R, ``Propagation data and prediction methods for the design of Earth-space
telecommunication systems,'' Recommendation ITU-R P.618--13, Geneva, Switzerland, 2017.
[SINR definition $\gamma = C/(N+I)$, rain attenuation, feeder-link margins]

\bibitem{goldsmith2005}
A.~Goldsmith, \emph{Wireless Communications}, Cambridge University Press, 2005.
[Multi-user SINR in linear domain; noise and interference are additive in linear scale, not dB]

\bibitem{3gpp_ntn}
3GPP, ``Study on New Radio (NR) and its evolution for Non-Terrestrial Networks,''
Technical Report 3GPP TR 38.821, Release 16, April 2020.
[Multi-beam HAPS interference model: $\mathrm{SINR}_k = P_{\mathrm{serving}} \cdot G_k / (N_0 + \sum P_i \cdot G_{i,k})$]

% === Implementation Libraries ===
\bibitem{itur_package}
I.~Portillo-Ramirez, ``itur: Python implementation of ITU-R P recommendations,''
GitHub, 2024. Available: \url{https://github.com/iportillo/itur}
[Bundled ITU-R models for rain, fog, gas, scintillation]

% === Deep Reinforcement Learning (Phase 2) ===
\bibitem{mnih2015dqn}
V.~Mnih, K.~Kavukcuoglu D.~Silver, A.~A.~Rusu, J.~Veness, G.~Bellemare,
A.~Graves, M.~Riedmüller, A.~K.~Fidjeland, G.~Ostrovski, S.~Petersen,
C.~Beattie, A.~Sadik, I.~Antonoglou, H.~King, D.~Kumaran, D.~Wierstra,
S.~Legg, and D.~Hassabis, ``Human-level control through deep reinforcement
learning,'' \emph{Nature}, vol.~529, pp.~529--533, 2015.
[Foundational DQN algorithm: experience replay, target networks, Q-learning]

\bibitem{sutton2018rl}
R.~S.~Sutton and A.~G.~Barto, \emph{Reinforcement Learning: An Introduction},
2nd~ed. Cambridge, MA: MIT Press, 2018.
[MDP theory, value functions, policy gradient methods; Chapter 6: Q-learning]

\bibitem{schulman2017ppo}
J.~Schulman, F.~Wolski, P.~Dhariwal, A.~Radford, and O.~Klimov,
``Proximal policy optimization algorithms,'' arXiv preprint arXiv:1707.06347, 2017.
[Policy-gradient alternative to DQN; higher sample complexity for discrete actions]

\bibitem{goldsmith2005wireless}

\bibitem{mnih2015human}
V.~Mnih et~al., ``Human-level control through deep reinforcement learning,''
\textit{Nature}, vol.~529, no.~7587, pp.~529--533, Feb.~2015.

\bibitem{jain1984quantitative}
R.~K.~Jain et~al., ``A quantitative measure of fairness and discrimination for resource allocation in shared computer systems,''
\textit{arXiv preprint arXiv:cs/9809099}, 1984.

\bibitem{sutton2018reinforcement}
R.~S.~Sutton and A.~G.~Barto, \textit{Reinforcement Learning: An Introduction},
2nd~ed. MIT Press, 2018.
A.~Goldsmith, \emph{Wireless Communications}, 2nd~ed. Cambridge University Press, 2005.
[Chapter 3: SINR computation in linear domain, not dB; multi-user interference]

\end{thebibliography}
